\chapter{System Design and Pipeline Architecture}
\label{ch:design}

The pipeline consists of thirteen numbered, single-purpose scripts in
\texttt{v2-proj/}. Each script reads from upstream outputs and writes its own
results to \texttt{v2-proj/output/}. No script modifies the original source
data. All random seeds are fixed for reproducibility.
Figure~\ref{fig:pipeline_dag} shows the data flow.

\begin{figure}[H]
\centering
\usetikzlibrary{arrows.meta,positioning,fit,backgrounds}
\begin{tikzpicture}[
  node distance=0.55cm and 1.1cm,
  box/.style={rectangle, draw=ATUGreen, thick, rounded corners=3pt,
              fill=ATULightGreen!40, text width=3.8cm, align=center,
              font=\small, inner sep=4pt},
  data/.style={rectangle, draw=ATUNavy, thick, rounded corners=3pt,
               fill=ATUTeal!20, text width=3.2cm, align=center,
               font=\footnotesize\itshape, inner sep=3pt},
  arr/.style={-{Latex[length=2mm]}, thick, ATUGreen},
]

%--- Raw data
\node[data] (raw) {Raw data\\MIMIC-IV + eICU};

%--- Pipeline stages
\node[box, below=of raw]          (s01) {\texttt{01} Setup};
\node[box, below=of s01]          (s02) {\texttt{02} Cohort Labels};
\node[box, below=of s02]          (s03) {\texttt{03} Person-Hours};
\node[box, below=of s03]          (s07) {\texttt{07} Sample Size};
\node[box, below=of s07]          (s08) {\texttt{08} Primary Model};
\node[box, below=of s08]          (s09) {\texttt{09} Comparators};
\node[box, below=of s09]          (s10) {\texttt{10} Metrics Suite};
\node[box, below=of s10]          (s11) {\texttt{11} External Valid.};
\node[box, below=of s11]          (s12) {\texttt{12} Variance Decomp.};
\node[box, below=of s12]          (s13) {\texttt{13} Equity Analysis};

%--- Output column
\node[data, right=of s02]         (cohort)  {cohort\_\{A,B,C\}.parquet};
\node[data, right=of s03]         (ph)      {person\_hours\_\{A,B,C\}.parquet};
\node[data, right=of s08]         (models)  {model\_primary\_\{A,B,C\}.rds};
\node[data, right=of s09]         (preds)   {predictions\_comparators.parquet};
\node[data, right=of s10]         (metrics) {10\_metric\_results.csv};
\node[data, right=of s11]         (ext)     {11\_external\_validation.csv};
\node[data, right=of s12]         (decomp)  {12\_variance\_decomp.csv};

%--- Arrows (pipeline flow)
\draw[arr] (raw)  -- (s01);
\draw[arr] (s01)  -- (s02);
\draw[arr] (s02)  -- (s03);
\draw[arr] (s03)  -- (s07);
\draw[arr] (s07)  -- (s08);
\draw[arr] (s08)  -- (s09);
\draw[arr] (s09)  -- (s10);
\draw[arr] (s10)  -- (s11);
\draw[arr] (s11)  -- (s12);
\draw[arr] (s12)  -- (s13);

%--- Output arrows
\draw[arr] (s02)  -- (cohort);
\draw[arr] (s03)  -- (ph);
\draw[arr] (s08)  -- (models);
\draw[arr] (s09)  -- (preds);
\draw[arr] (s10)  -- (metrics);
\draw[arr] (s11)  -- (ext);
\draw[arr] (s12)  -- (decomp);

%--- Feed-forward arrows
\draw[arr] (ph.south) to[out=-90,in=180] (s08.west);
\draw[arr] (ph.south) to[out=-90,in=180] (s09.west);

\end{tikzpicture}
\caption{Pipeline data flow. Each box is a single script; italic nodes are
  output files. Raw MIMIC-IV and eICU data are read-only.}
\label{fig:pipeline_dag}
\end{figure}

\section{Configuration and Utilities}

\texttt{config.R} is the single source of truth for all pipeline settings:
data paths, cohort criteria, antibiotic name lists, MIMIC-IV item identifiers,
label variant definitions, split boundaries, prediction framing constants,
equity variable names, and Utility Score parameters. No other script hard-codes
a clinical constant or file path.

\texttt{utils.R} provides shared helper functions: a DuckDB-backed streaming
reader for large CSVs, Parquet reader/writer wrappers, datetime parsing,
last-observation-carried-forward (LOCF) imputation, and Schoenfeld test
wrappers.

The full pipeline can be reproduced on any MIMIC-IV installation by setting a
single environment variable (\texttt{MIMIC\_RESEARCH\_ROOT}).

\section{Stage 01: Setup Validation}

\texttt{01\_setup.Rmd} checks that all required MIMIC-IV source tables exist
and reports their sizes. The pipeline stops if any table is missing.

\section{Stage 02: Cohort Extraction and Labelling}

\texttt{02\_cohort\_labels.Rmd} builds all three cohorts in a single pass:

\begin{enumerate}
  \item \textbf{Base cohort:} Adult first ICU stays with LOS $\geq 4$\,h
  (65,078 stays). Equity variables (race, language, insurance) are attached.

  \item \textbf{Suspected infection:} Qualifying antibiotic-culture pairs
  within each variant's time windows. Surveillance cultures are excluded.

  \item \textbf{SOFA organ dysfunction:} SOFA components are computed from
  chartevents, labevents, inputevents, and outputevents. A SOFA rise
  $\geq 2$ within each variant's windows defines organ dysfunction.

  \item \textbf{Onset time:} The hour when both criteria (infection + SOFA rise)
  are simultaneously met.
\end{enumerate}

Outputs: \texttt{cohort\_A/B/C.parquet} and \texttt{label\_variance\_table.csv}.

\section{Stage 03: Person-Hour Table Construction}

\texttt{03\_person\_hours.Rmd} expands each stay into hourly rows. Large source
tables (chartevents: 3.5\,GB; labevents: 2.6\,GB) are processed via DuckDB
streaming, pre-filtered to cohort stays, to avoid loading everything into
memory. For each person-hour:

\begin{itemize}
  \item Vitals and labs are LOCF-imputed from the most recent measurement.
  \item Six-hour rolling slopes are computed for respiratory rate, MAP, and
  lactate.
  \item Vasopressor indicators are set from concurrent infusions.
  \item Missingness flags are set for each lab test.
  \item The competing-event code (0/1/2/3) is assigned based on the hour's
  relationship to onset, discharge, or death.
  \item The temporal split group is assigned from the stay's admission year.
  \item Treatment anchor timestamps (first antibiotic, first culture) are
  recorded.
\end{itemize}

Outputs: \texttt{person\_hours\_A/B/C.parquet}.

\section{Stage 04: Sample Size Calculation}

\texttt{04\_sample\_size.Rmd} applies the Riley et al.\
\cite{riley2019sample} criteria to the training portion of each
person-hour table (2008--2016). Sufficiency is evaluated per label variant.

\section{Stage 05: Primary Model}

\texttt{05\_primary\_model.Rmd} fits the multinomial discrete-time
competing-risks model (Equation~\ref{eq:primary_model}) for each label variant
on the temporal training split. Baseline hazard splines use 5 knots. Standard
errors are cluster-robust (grouped by stay). A random patient-level 75/25 split
is also fitted to support H4.

Outputs: \texttt{model\_primary\_A/B/C.rds}, coefficient tables, and
\texttt{predictions\_primary\_A/B/C.parquet}.

\section{Stage 06: Comparator Models}

\texttt{06\_comparators.Rmd} fits the XGBoost GBT model, the static Cox model
(with Schoenfeld testing), and computes NEWS2, qSOFA, and SIRS at each
person-hour. GBT models are fitted under both temporal and random splits. GBT
feature importance is exported for each variant.

Outputs: model objects, predictions, feature importance tables, and Schoenfeld
test results.

\section{Stage 07: Metrics Suite}

\texttt{07\_metrics\_suite.Rmd} evaluates all models on all splits (temporal,
random, treatment-anchored) using the full metric set: AUROC, AUPRC, calibration
intercept and slope, PhysioNet Utility Score (via Python module
\texttt{utility\_score.py}), alert burden, lead time, and decision-curve net
benefit.

Output: \texttt{10\_metric\_results.csv}.

\section{Stage 08: External Validation on eICU-CRD}

\texttt{08\_external\_validation.Rmd} applies the frozen MIMIC-IV models
(primary and NEWS2) to eICU-CRD without any retraining. AUROC and AUPRC are
reported for each variant and model.

\section{Stage 09: Variance Decomposition}

\texttt{09\_variance\_decomposition.Rmd} computes the pre-registered variance
decomposition by varying one factor at a time while holding all others fixed:

\begin{itemize}
  \item \textbf{Label spread (H1):} AUROC range across Variants A/B/C.
  \item \textbf{Model-class spread:} AUROC range across all models.
  \item \textbf{Split optimism (H4):} Random vs.\ temporal split comparison.
  \item \textbf{Metric divergence (H5):} AUROC vs.\ Utility Score decline.
\end{itemize}

Outputs: decomposition tables, hypothesis verdicts, and comparison plots.

\section{Stage 10: Equity Analysis}

\texttt{10\_equity\_analysis.Rmd} evaluates model performance and label
sensitivity within each stratum of four equity variables: race, ethnicity,
language, and insurance. The combination of performance disparities and
label-sensitivity variation is the novel finding.

\section{Reproducibility}

The pipeline is orchestrated by \texttt{run\_pipeline.sh}. All intermediate
outputs are Parquet (compressed, columnar). Small result tables ($<$1,000 rows)
are CSV for readability. The full pipeline runs in approximately six hours on
a machine with 64\,GB RAM.
